% Section: Discussion — P8 commit 8.
% Cross-references existing committed numbers; no new claims that
% require integrity gates. Limitations + future work + relation to
% the broader QML benchmarking literature.

\section{Discussion}
\label{sec:discussion}

The pre-registered $H_1$ regime-dependent quantum advantage
hypothesis is FALSIFIED across all eight non-fragile verdicts
(Table~\ref{tab:master-verdicts}). The LTC decomposition
(Section~\ref{sec:forecaster-decomp}) localizes the mechanism in
the quantum circuit; the H3 cross-tabulation
(Section~\ref{sec:mechanism}) finds a candidate
expressivity-driven trend, underpowered at $n=9$. Together these
constitute a rigorous mechanistic null on a falsifiable quantum-ML
hypothesis --- the kind of result the Bowles--Schuld program
calls for~\cite{bowles2024better}.

\subsection{Limitations}
\label{sec:discussion-limitations}

\textbf{Sample size.} The forecaster task and the HPO-best
sensitivity point run at $n=9$; the combined ODE+PDE solver
verdict at $n=24$ (Section~\ref{sec:solver-verdicts}) is the
largest. A larger ladder would tighten the CIs but, given the
consistency across eight verdicts, is unlikely to shift the
qualitative outcome.

\textbf{Deferred systems.} Two pre-registered systems were
gate-tested but deferred for compute (amendment A11):
\begin{itemize}
  \item \textbf{Kuramoto (12-dim)} -- the per-component scalar
        circuit scales linearly with state dimension; a 12D
        system would cost $\approx 7$\,hr per cell. A single-cell
        smoke test confirmed the dispatch works.
  \item \textbf{KdV ($u_t + 6 u u_x + u_{xxx} = 0$)} -- the
        \texttt{jacrev$^3$} triple-nested autodiff through the
        QNode gate-tested
        (\texttt{results/p7\_8\_kdv\_gate/gate\_result.json})
        \emph{passes}: finite $u_{xxx}$ at JIT compile time
        $\approx 4.6$\,s and per-point cost $\approx 0.43\times$
        the second-order baseline via XLA fusion. Integrated
        training cost remains prohibitive at $\approx 12$\,s per
        gradient step.
\end{itemize}

\textbf{Matched-capacity bound.} Forecaster comparisons hold
parameter count within a factor of two
(Section~\ref{sec:methods-baselines}). The cPINN solver baseline
violates this by $3.3\times$ (amendment A4) in the classical
baseline's favor, and is therefore conservative for the
FALSIFIED outcome.

\textbf{Hyperparameter optimization.} P7.6 ran symmetric HPO on
the solver task at three anchor cells per family (LV s2,
\vdp{} s1, Lorenz s2), three learning rates $\times$ two
step budgets; its HPO-best $n=9$ verdict (row 4 of
Table~\ref{tab:master-verdicts}) is the methodologically
strongest sensitivity in the paper. The forecaster task uses
default-Adam; an HPO-symmetric forecaster sensitivity is a
natural revision-stage extension.

\textbf{Simulator only.} All quantum-circuit evaluations use
\texttt{default.qubit} in PennyLane, which the
pre-registration~\cite{schuld2022right} explicitly disclaims
hardware execution for. A small-scale hardware run (e.g.,
data-reuploading at $n_q = 4$) would be a natural revision-stage
validation but does not address the falsified hypothesis.

\textbf{Single ansatz per family.} Each ansatz uses one layered
configuration $(n_q, L) = (4, 1)$ at the default budget.
Architectural variants (deeper reuploading at $L = 5$,
alternative entanglement patterns) are deferred to the follow-up
paper.

\textbf{Multiple-comparison family.} The eleven verdicts of
Table~\ref{tab:master-verdicts} form a single inferential family
that the per-row $\alpha = 0.05$ CI machinery does not control.
Under strict Holm--Bonferroni ($\alpha_{\text{family}} = 0.05$,
$m = 11$), both PRIMARY headline verdicts retain FALSIFIED
designation and the eight ancillary verdicts already crossed
zero, so the correction is conservative; it is disclosed for
transparency (Section~\ref{sec:methods-verdict} and the
Table~\ref{tab:master-verdicts} footnote).

\textbf{Latent-bug remediation.} The post-peer-review audit
(Amendments A20--A22) surfaced three issues that left the
committed headline numbers intact: the \texttt{te\_qpinn\_fnn}
readout-convention divergence (A20, requiring cell re-runs), the
\texttt{brickwall} connectivity reduction at $n_q = 3$ (A21,
disclosure-only strengthening of A18), and a 2D PDE hard-IC
docstring omission (A22, latent; implementation always correct).
The post-A20 re-runs are bundled into the Phase~C Anvil sweep
(Appendix~A).

\subsection{Future work: hardening quantum PINN training}
\label{sec:discussion-future}

The mechanism analysis (Section~\ref{sec:mechanism}) identifies
three interventions from the quantum-ML literature as plausibly
relevant to the regime-contrast inversion:

\begin{enumerate}
  \item \textbf{Quantum Natural
        Gradient}~\cite{stokes2020quantum}. Adam treats the
        variational parameter manifold as Euclidean; QNG
        preconditions with the Quantum Fisher Information matrix,
        improving convergence and loss floors. The H3 KL-to-Haar
        trend $\rho = +0.518$ (Section~\ref{sec:mechanism-h3})
        is consistent with the expressivity-helps reading QNG
        targets.

  \item \textbf{Causal training
        weighting}~\cite{wang2022causality}. The PINN residual
        loss weights all collocation points equally, ignoring
        temporal ordering; an inverse-time-order schedule that
        grows weights only after upstream residuals decay
        improves training on stiff dynamics, directly the \vdp{}
        failure mode (Section~\ref{sec:forecaster-cells}).

  \item \textbf{Higher data-reuploading
        depth}~\cite{schuld2021effect}. H3 found Fourier
        $K_{\max} = 3$ at $L = 1$ across all four ansätze --- the
        ceiling at the current configuration. The
        Schuld--Sweke--Meyer relation $K_{\max} = L \cdot n$
        predicts $K_{\max} = 15$ at $L = 5$. If the smooth-bin
        inductive bias the LTC decomposition isolated
        ($\Delta_\text{liquid} = +0.12$,
        Section~\ref{sec:forecaster-decomp}) scales with
        $K_{\max}$, deeper reuploading should recover the $H_1$
        direction.
\end{enumerate}

The negative finding and its mechanistic substructure
(substrate-dependent $\tau$ signature, KL-to-Haar trend) seed
three follow-up papers, listed by experimental readiness:

\begin{enumerate}
  \item \textbf{The training-hardening paper.} The three
        interventions above (QNG~\cite{stokes2020quantum}, causal
        weighting~\cite{wang2022causality},
        higher-$K_{\max}$~\cite{schuld2021effect}) are tested
        jointly on the same hardness ladder --- plus the deferred
        Kuramoto and KdV systems --- at P7.6's HPO compute budget,
        with the same
        pre-registration and verdict machinery (locked in
        \texttt{.planning/P7\_7\_HARDENING\_PREREG.md}). If any
        subset recovers $H_1$ in the predicted direction the
        paper reports it; otherwise it strengthens the present
        FALSIFIED verdict from one to four sensitivity points.
        Scope: $\sim 60$ cells, $\sim 18$~CPU-hours on Anvil.

  \item \textbf{The $\tau$-substrate mechanism paper.} The
        substrate-dependent $\tau$ signature
        (Figure~\ref{fig:tau-substrate}; $\Delta_{\mathrm{liquid}|\,
        \mathrm{classical}} \approx +0.12$ vs.\
        $\Delta_{\mathrm{liquid}|\,\mathrm{quantum}} \approx -0.33$,
        $|\textrm{gap}| = 0.45$) is the strongest mechanism finding
        here and is explained by neither substrate alone. The
        follow-up isolates the LTC $\tau$ machinery across substrate
        widths and reuploading depths, seeking a closed-form
        prediction for the sign of $\Delta_\tau$ in the spirit of
        Hasani et al.~\cite{hasani2022closed}. Scope: $\sim 36$
        cells.

  \item \textbf{The domain-breadth paper (ChemE / BioE / PSE).}
        Extends the ladder from 8 dynamical-system archetypes to 12
        process-engineering ODEs / PDEs spanning chemical-reactor
        (Van~de~Vusse, Saponification), biochemical (Monod,
        Yeast / Andrews, Contois), and process-systems
        (Fisher--KPP, CDR low/high-Pe, fixed-bed, ADM1) regimes,
        testing whether the FALSIFIED H1 generalizes to the
        domain-specific ladder a practitioner cares about. Scope:
        $\sim 12$ systems $\times$ 4 families $\times$ 3 seeds
        $= 144$ cells, $\sim 36$~CPU-hours.
\end{enumerate}

All three inherit this paper's reproducibility infrastructure:
pre-registration template, matched-capacity / paired-bootstrap /
underfit-skyline-guard verdict machinery, and the
\texttt{verify\_paper\_integrity} mechanical gate.

\subsection{Positioning relative to the field}
\label{sec:discussion-position}

\textbf{(1) The FALSIFIED outcome is not a uniqueness claim.}
The specific regime-dependent advantage formulated by
Schuld--Killoran~\cite{schuld2022right} does not materialize on
this ladder at matched-capacity training. This neither shows that
all quantum advantages on PINN-style tasks fail nor that the
Schuld--Fourier hypothesis is wrong in general --- the
liquid-isolated decomposition ($\Delta = +0.12$) is the one
verdict whose point estimate matches the original direction, which
the higher-$K_{\max}$ intervention will probe.

\textbf{(2) The mechanism attribution is the scientific
content.} A FALSIFIED pre-registered hypothesis is
publishable~\cite{bowles2024better}, but its weight depends on
what the falsification \emph{shows}.
Sections~\ref{sec:forecaster-decomp} and \ref{sec:mechanism}
attribute the regime-contrast inversion to the quantum circuit ---
not the liquid-$\tau$ machinery, not baseline-tuning artifacts ---
a contribution to understanding when matched-capacity quantum-PINN
training fails.

\textbf{(3) The positive sub-findings are independent of the
aggregate verdict.} \texttt{te\_qpinn\_qnn} achieves a
$24\times$-tighter seed standard deviation and $2.5\times$ lower
mean relative-$L^2$ than the matched classical PINN on \fhn{}
(Section~\ref{sec:solver-fhn}); \texttt{rf\_qrc} achieves the
strongest cell on Lorenz s0 and s2
(Section~\ref{sec:forecaster-cells}). Both occur at zero classical
parameters in the quantum-cell stack --- regime-specific
structural advantages the falsification does not contradict, of
the kind the \texttt{te\_qpinn\_qnn}~\cite{berger2025teqpinn} and
\texttt{rf\_qrc}~\cite{ahmed2024rfqrc} original works target, here
seen in an independent benchmark.

\textbf{(4) Comparison to classical operator-learning competitors.}
Our matched-classical baselines are MLP-PINNs, mirroring the
PINN-vs-PINN framing of the quantum-PINN literature itself.
Benchmarking against operator-learning state-of-the-art ---
DeepONet~\cite{lu2021deeponet} and the Fourier Neural
Operator~\cite{li2021fourier} --- would collapse the PINN
inductive-bias question (trial-solution structure) into a
representational-capacity question (data scale and architecture);
the follow-up paper is the appropriate venue for such baselines,
paired with the corresponding generalization
analysis~\cite{caro2022generalization}.

\subsection{Reproducibility and audit}
\label{sec:discussion-reproducibility}

Every reported number is gated mechanically by
\texttt{scripts/verify\_paper\_integrity.py}, which matches
committed JSON files against the paper's text; a reviewer running
it after cloning the repository
(\url{https://github.com/shawngibford/qlnn}) receives a
pass/fail summary on every numerical claim, including the LTC
decomposition's algebraic-identity check
(Section~\ref{sec:forecaster-decomp}). Twenty-two amendments
(A1--A22) are documented in \texttt{PRE\_REG\_AMENDMENT.md} with
per-amendment empirical findings: A15--A19 (2026-05-28 internal
audit) propose budget-parity, ansatz-aliasing, and
sub-experiment refinements; A20--A22, filed the same day after a
five-reviewer pass, record a paired QPINN readout divergence, a
brickwall connectivity diagnosis, and a latent docstring defect.
The verdict refresh under the A15--A22 configuration is deferred
to a supplementary compute campaign.

The full reproduction pipeline
(\texttt{scripts/reproduce\_paper.sh}) regenerates every cell
from scratch on a single MacBook Pro M1 in $\approx 8$\,hours,
parallelizable across sweep scripts; the 36 P7.11 non-liquid
quantum cells reproduce in under 30 minutes.
